\section{The Custom three.js Simulation Engine}
\label{sec:webengine}

The research pipeline described in the preceding chapters produces its
results through batch experiments: a simulator steps through planning
cycles, records outcomes, and reports statistics after the fact. Such
experiments establish what the pipeline \emph{does}, but they are poorly
suited to exploring how it \emph{behaves}: how pruned topologies respond to
faults, how trust estimates track hidden reliability, or how competing
assignment policies diverge when driven by identical evidence. To make
these dynamics inspectable, we built an interactive browser console that
drives the \emph{real} research pipeline, comprising constrained QUBO edge
pruning, trust-aware GNN scoring, greedy task assignment, and trust model
updates, through a Python backend, while the browser itself owns only
kinematics and rendering. The console is not a visualization bolted onto
recorded data; every cycle shown on screen is produced by the same code
that generated the results in this report. We regard the engine as a work
in progress, and this chapter describes its current state, including
explicitly marked placeholders for material that remains to be captured.

\fillme{screenshot: orbit view of the fleet console with link arcs visible}

\subsection{Overview and Design Principles}

The console is a single-page web application written in TypeScript with
React~18, rendered through three.js via react-three-fiber, with
application state held in a Zustand store and the project bundled by Vite.
The backend is a deliberately dependency-free Python standard-library HTTP
server, chosen so that the console runs anywhere the research code runs,
without a package installation step beyond what the pipeline already needs.
Three principles guided the design.

First, \emph{the browser never plans}. All optimization, scoring, and
outcome sampling happen in the Python session server; the client computes
rover motion and draws the scene. This separation guarantees that what the
user sees is the pipeline under test, not a reimplementation that might
quietly diverge from it. Second, \emph{one source of truth for
configuration}. A shared configuration file, read by both the TypeScript
client and the Python server, defines the arena geometry, vehicle
kinematic limits, every tunable parameter with its scope and default, the
calibrated presets, and the comparison catalogue with its published
reference numbers. The two processes cannot disagree about, say, the
rover's turning limit, because both read the same numbers. Third,
\emph{honest comparison by construction}. The console implements the
paired-comparison methodology of the audit rather than a looser visual
side-by-side, and it discloses the baselines it does not offer rather than
hiding them.

\subsection{Two Processes, One Command}

A single development script starts both processes: the session server and
the Vite development server. At runtime the two communicate over a small
JSON API. The client polls \code{GET /api/health} at startup and
periodically thereafter; the response reports which optional solvers are
installed (neal, tabu, SQA, MILP), so the comparison menu only offers
baselines that can actually run on the current machine. Each planning
cycle is requested with \code{POST /api/step}, whose response carries the
kept-edge topology, assignment, per-rover scores and trust values, outcome
realizations, and wall-time measurements for every active baseline.

\begin{table}[t]
\centering
\caption{HTTP API of the session server. All validation failures return a
structured error body; no response ever contains a NaN literal.}
\label{tab:api}
\begin{tabular}{@{}llp{5.6cm}@{}}
\toprule
Method and path & Error & Meaning \\
\midrule
\code{GET /api/health} & 200 & Reports installed optional solvers and session load. \\
\code{POST /api/step} & 400 & Body fails validation (malformed JSON, body over 32~KiB, non-finite or out-of-arena positions, parameters out of range, more than 8 faults). \\
\code{POST /api/step} & 429 & All 16 session slots are occupied and none is evictable. \\
\code{POST /api/step} & 503 & The solve semaphore is saturated (4 concurrent solves); retry later. \\
\bottomrule
\end{tabular}
\end{table}

Because the console is interactive, the server is hardened against inputs
a batch script would never produce. Request bodies are capped at 32~KiB;
rover positions must be finite and inside the arena; every tunable
parameter is range-checked against the shared configuration; and at most
8~faults may be injected per cycle. Concurrency is bounded by a semaphore
admitting four simultaneous solves, with \code{503} returned when the
server is busy, so a user dragging a slider rapidly cannot queue an
unbounded solve backlog. The server hosts at most 16~sessions, each with a
600~s time-to-live and busy-window eviction, returning \code{429} when
full. Finally, every outgoing payload passes through a sanitizer, so NaN
or infinite values, which arise legitimately inside numerical code during
degenerate configurations, can never leak into a JSON response, where they
would be unparseable. Table~\ref{tab:api} summarizes the surface.

\begin{figure}[t]
\centering
\resizebox{0.95\textwidth}{!}{%
\begin{tikzpicture}[
  font=\small,
  box/.style={draw, rounded corners=2pt, align=center, minimum height=0.85cm, inner sep=4pt},
  layer/.style={draw, dashed, rounded corners=4pt, inner sep=6pt},
  arr/.style={-{Latex[length=2mm]}, thick},
  every label/.style={font=\footnotesize\itshape}
]
% Browser engine boxes
\node[box] (core) {Engine core\\ \footnotesize fixed-step kinematics, spawn, delivery};
\node[box, right=0.5cm of core] (render) {Render layer\\ \footnotesize instanced three.js scene};
\node[box, below=0.45cm of core] (state) {State store\\ \footnotesize Zustand, 240-cycle ring};
\node[box, below=0.45cm of render] (ui) {UI panels\\ \footnotesize mission, tuning, roster, minimap};

\begin{scope}[on background layer]
  \node[layer, fit=(core)(render)(state)(ui),
        label={[anchor=north west]north west:Browser (React, three.js)}] (browser) {};
\end{scope}

% Server side
\node[box, right=2.6cm of render] (server) {Session server\\ \footnotesize stdlib HTTP, 16 sessions, 4-solve semaphore};
\node[box, below=0.9cm of server] (pipeline) {Research pipeline\\ \footnotesize QUBO pruning, GNN scoring,\\ \footnotesize trust updates, outcomes};

% Arrows
\draw[arr] (core) -- node[above, font=\footnotesize]{\code{POST /api/step}} (server);
\draw[arr] (server) -- node[below, font=\footnotesize]{topology, scores, outcomes} (core);
\draw[arr] (server) -- node[right, font=\footnotesize]{one cycle} (pipeline);
\draw[arr] (pipeline) -- node[left, font=\footnotesize]{results} (server);
\draw[arr] (core) -- (render);
\draw[arr] (core) -- (state);
\draw[arr] (state) -- (ui);
\draw[arr] (ui) -- node[below, font=\footnotesize]{tuning, faults} (state);
\end{tikzpicture}}
\caption{Console architecture. The browser owns kinematics, rendering,
state, and interaction; the session server owns planning. A single shared
configuration (not shown) is read by both sides.}
\label{fig:webarch}
\end{figure}

Figure~\ref{fig:webarch} shows the division of labor. The browser engine
splits into four cooperating layers: a pure, testable core implementing
kinematics, spawning, and delivery; a render layer built on instanced
three.js geometry; a central state store; and the UI panels. Only the core
talks to the network.

\subsection{The Shared Configuration}

The shared configuration is the contract between client and server. It
fixes the physical world: a circular arena of radius 6~m, rovers with body
radius 0.62~m, wheel radius 0.16~m, and track 0.76~m, limited to a maximum
wheel speed of 1.25~m/s, acceleration 0.9~m/s$^2$, braking 1.8~m/s$^2$,
and yaw rate 1.8~rad/s. On top of this it defines 15~tunable parameters,
each annotated with a \emph{scope} that the UI exposes to the user:
session-scoped parameters restart the mission when changed (for example
fleet size), planner-scoped parameters take effect on the next planning
cycle (for example the QUBO budget or solve interval), and client-scoped
parameters apply instantly (for example camera mode or render quality).

Seven calibrated presets ship with the console
(Table~\ref{tab:presets}), spanning benign field operations through a
deliberately hostile communications regime. Each preset carries a measured
per-cycle planning time on our reference machine, ranging from 140~ms for
the lightest configuration to 1420~ms for the largest fleet, so users can
see the real cost of the solver settings they choose rather than
discovering it as interface lag. The same configuration also holds the
comparison catalogue: for every baseline the console can run, it stores
the published reference numbers from the audit, which the mission strip
displays alongside live results.

\begin{table}[t]
\centering
\caption{Calibrated mission presets. Planning times are measured per-cycle
medians on our reference machine; where an exact value is not listed here
we give the span across all presets, 140~ms to 1420~ms.}
\label{tab:presets}
\begin{tabular}{@{}lll@{}}
\toprule
Preset & Purpose & Cycle time (ms) \\
\midrule
field   & Benign outdoor survey, default mission            & 140 to 300 \\
rapid   & Fast replanning, solve every cycle                & 140 to 400 \\
relay   & Stretched topology stressing connectivity         & 300 to 700 \\
swarm   & Many small rovers, dense graph                    & 400 to 1000 \\
fleet16 & Sixteen rovers, upper end of the live fleet       & 900 to 1420 \\
hostile & High packet loss, frequent faults, drift          & 300 to 900 \\
tight   & Small QUBO budget, aggressive pruning             & 140 to 500 \\
\bottomrule
\end{tabular}
\end{table}

\subsection{Honest Comparison by Construction}

A demo that lets the user pit our pipeline against alternatives invites an
obvious failure mode: comparisons that look dramatic but are confounded.
The console addresses this structurally. Assignment baselines
(trust-only, random, and no-graph) do not run independent missions; they
run as \emph{paired follower simulations} that share the primary mission's
kept-edge topology, so they receive identical situational evidence and
differ in assignment policy only. Faults injected from the UI, whether a
reliability drop, battery drain, or a packet loss spike, are applied
identically to the primary and all followers, preserving the pairing.
Solver baselines (bit-level SA, neal, tabu, and SQA) run on the
\emph{exact same $Q$ matrix} the production annealer just solved, seeded
with the primary's previous solution as a warm start, and the console
reports for each of them the energy gap, feasibility, matched-or-better
rate, and wall time.

The comparison menu also embodies the audit's honesty policy in a
deliberately uncomfortable way: it excludes the two baselines that
\emph{beat} our pipeline, the exact MILP solve and the simple threshold
pruner, and says so in a footnote beneath the menu rather than concealing
their existence. Users exploring the console see the same negative results
the paper reports.

\fillme{screenshot: mission strip during a live solver comparison}

\subsection{The Engine Core}

Motion and delivery are computed by a custom fixed-timestep engine written
in pure TypeScript. Simulation advances at 1/120~s steps through an
accumulator, with rendering interpolating between the two most recent
states, so visual smoothness is decoupled from both the integration rate
and the planner cadence. All simulation state lives in typed arrays rather
than per-object JavaScript properties, which keeps memory layout flat,
avoids per-frame allocation, and makes the hot loops friendly to the JIT.

\paragraph{Kinematics.}
Rovers are differential-drive vehicles. Rather than the common Euler
approximation ($x \mathrel{+}= v\cos\theta\,\Delta t$ and so on), we
integrate each step in exact constant-twist form: over a step the rover
travels a circular arc determined by its linear and angular velocity, and
the pose update is closed-form in the arc radius $v/\omega$. The naive
formula divides by the yaw rate $\omega$, so we branch on
$|\omega| < \varepsilon$ and evaluate the singular case through a
$\operatorname{sinc}$ expansion, making the integrator singularity-free at
zero yaw rate. Wheel-speed commands that exceed the 1.25~m/s limit are
saturated by scaling \emph{both} wheels by the same factor, which
preserves curvature: the rover follows exactly the same arc, only slower,
rather than veering off the intended path. An asymmetric rate limiter
enforces the 0.9~m/s$^2$ acceleration and 1.8~m/s$^2$ braking limits, so
rovers stop faster than they start, a small asymmetry that reads as
mechanically plausible on screen.

\paragraph{Steering and collision avoidance.}
The steering controller sums several influences: seek toward the current
goal, quadratic repulsion from neighboring rovers, a keep-right rule for
head-on encounters, a soft potential at the arena wall, and, for missions
that want it, an optional tether pulling a rover toward its nearest
network neighbor so the fleet does not silently fragment. A speed profile
slows the rover into turns and as it approaches an arrival. Two hard
guarantees sit beneath all of this. First, a \emph{collision invariant}:
no pair of rovers may ever come closer than two body radii. Any
integration step that would violate this is rejected, and the rover
rotates in place instead, so interpenetration is impossible by
construction regardless of how aggressive the steering sum becomes.
Second, \emph{stuck detection}: a rover whose progress stalls executes
alternating-side detours, escaping the symmetric deadlocks that pure
repulsion fields are prone to.

\paragraph{Spawning and task placement.}
Missions begin from a deterministic spawn layout produced by dart throwing
with a hexagonal lattice fallback, so a given fleet size and seed always
yield the same starting configuration, keeping paired comparisons fair.
Task sites are derived from a stable hash of the task identifier rather
than from a stream of random draws, which means the same task id lands in
the same place in the primary and in every follower simulation.

\paragraph{Delivery model.}
The browser reproduces the research harness's delivery semantics: message
success follows a widest-path (maximum bottleneck) tree from each rover to
the root rover, so the console's on-screen success statistics are the same
quantity the offline simulator measures, not a decorative substitute. Task
beacons carry 14~s timeouts, after which an unclaimed task expires and is
recorded as a failure.

\subsection{The Render Layer}

Interactivity at sixteen rovers with hundreds of candidate links requires
discipline about draw calls, so the render layer is \emph{instanced
everything}: one draw call per system regardless of fleet size.

All radio links, candidate, kept, breaking, and widest-path, are drawn by
a single shader-drawn ribbon mesh. Each link is a 24-segment quadratic arc
whose style is selected per instance: candidate links are faint, kept
links solid, breaking links flicker at 14~Hz to draw the eye, and the
widest-path tree is highlighted with animated packet flow moving toward
the root, giving an immediate read on where the network's bottlenecks
currently are. One additional instanced draw renders all ground decals:
contact shadows, communication-range rings, the selection ring, fault
glows, and beacon pads. Task beacons are a further instanced mesh.

The fleet itself is drawn as one instanced mesh per body part plus wheel
groups. The default rover model is procedural, with dimensions taken from
the physics simulator's rover description: a chassis of $0.9 \times 0.65
\times 0.25$~m on 0.16~m wheels, so the browser robot matches the simulated
robot rather than an artist's impression. A custom vertex-attribute PBR
technique bakes color, roughness, and metalness into vertex data,
collapsing what would conventionally be a multi-material, multi-draw model
into a single draw call for the entire body. Users can also import their
own GLB rover models; imports are validated against 24~MB and
250{,}000-triangle limits, wheels are detected by node name, and the model
is normalized to the standard footprint so a user-supplied robot obeys the
same kinematics as the default.

The camera offers orbit, top-down, follow, and chase modes with smooth
glides between them, switched from the keyboard (keys 1 to 4; space pauses
the mission). The view frustum's ground footprint is projected into an
RTS-style north-up minimap, so the user can see both what the main camera
covers and where every rover sits. Three quality tiers adapt rendering
cost to the hardware: high enables MSAA, ambient occlusion, and bloom;
balanced uses SMAA with bloom; low disables postprocessing entirely. A
performance monitor demotes the tier automatically when frame times sag,
and promotes it again when headroom returns.

\fillme{performance capture: fps and draw calls at 16 rovers on the high quality tier}

\subsection{State, Planner Loop, and User Interface}

Application state lives in a single Zustand store holding a 240-cycle
history ring, which backs the rolling charts without unbounded memory
growth. The planner loop keeps exactly one request in flight at an
adaptive cadence, 1.6~s by default; when the server reports congestion it
backs off exponentially to 8~s rather than piling requests onto a busy
backend. Independently, telemetry is copied into the store at 5~Hz so the
HUD stays fluid even while the planner is mid-cycle.

The interface comprises a top bar; a mission panel with tabs for mission
presets, comparisons and fault injection, tuning sliders annotated with
their scope notes and graph-theory advisories (for example warning when a
degree cap threatens connectivity), and view, quality, and model-import
controls; a mission strip showing live statistics, including bootstrap
confidence intervals over paired cycles, rolling sparklines, and
plain-language significance wording in place of raw $p$-values; a
per-rover roster whose heading-up radar canvases plot each rover's
estimated trust against the hidden ground-truth reliability, exposing how
fast the trust model converges after a fault; and the north-up minimap.
Faults injected from the panel, reliability drop, battery drain, or packet
loss spike, take effect in the primary and every follower simulation at
the same cycle, preserving the paired design.

\fillme{link or QR code to a demo video}

\subsection{Planned Engine Extensions}

The engine is functional but incomplete, and we close with the three
extensions we consider most valuable, in priority order.

\paragraph{WebSocket streaming for mid-cycle telemetry.}
The current poll-per-cycle design hides everything that happens inside a
solve. A WebSocket channel would stream intermediate annealer energies,
trust updates, and assignment proposals as they occur, turning the console
into an instrument for observing solver dynamics rather than only their
endpoints.

\paragraph{Session persistence and mission replay.}
Sessions currently evaporate with their 600~s TTL. Recording the full
request and response stream would support resuming interrupted sessions
and replaying missions deterministically, which would also simplify
regression testing of the client against recorded server behavior.

\paragraph{Motion command publishing.}
At present the physics simulator and the browser engine run separate world
models with matched parameters. Publishing motion commands from one engine
to the other, so both drive the same rovers, would remove the residual
risk of kinematic divergence and enable mixed-reality demonstrations in
which simulated and browser-rendered fleets interleave.
